\documentclass[11pt,a4paper,twoside,pdf]{article}

% Codificación y tipografía
\usepackage[T1]{fontenc}
\usepackage[utf8]{inputenc}
\usepackage{mathpazo}
\usepackage{braket}

% Matemáticas
\usepackage{amsmath, amssymb, amsfonts}
\usepackage{bm}
\usepackage[mathscr]{euscript}

% Figuras
\usepackage{graphicx}
\usepackage{float}
\usepackage{subcaption}

% Tablas y utilidades
\usepackage{multirow}
\usepackage{rotating}
\usepackage{array}
\usepackage{booktabs}
\usepackage{xspace}

% Símbolos adicionales
\usepackage{eurosym}
\usepackage{bbding}
\usepackage{pifont}

% Otros
\usepackage{soul}
\usepackage{color}
\usepackage{latexsym}
\usepackage{xtab}

% Hipervínculos
\usepackage[colorlinks=true,urlcolor=blue,linkcolor=blue,citecolor=blue]{hyperref}

% Numeración de ecuaciones
\numberwithin{equation}{section}

% Márgenes
\usepackage[top=2.88cm,bottom=2.97cm,left=2.95cm,right=2.95cm]{geometry}

% Idioma
\usepackage[spanish,es-nodecimaldot,es-tabla,es-lcroman,es-nosectiondot,es-noindentfirst]{babel}

% Cabeceras
\usepackage{fancyhdr}

% Teoremas
\newtheorem{theorem}{Teorema}[section]
\newtheorem{corollary}{Corolario}[theorem]
\newtheorem{lemma}[theorem]{Lema}

% Comandos propios
\newcommand{\dis}{\displaystyle}

\linespread{1.05}

\begin{document}

% Portada %%%%%%%%%%%%%%%%%%%%%%%%%%%%%%%%%%%%%%%%%%%%%%%%%%%%%%%%%%%%%%%%%%%%%%

\pagestyle{empty}

\noindent
\begin{tabular}{r}

\includegraphics[width=8.8cm]{escudoUGRmonocromo.png} \\[-1.8ex]
\hspace{31mm}\vspace{-8mm}
\begin{tabular}{c}
\hline\\[-1ex]\hskip-2mm
{\bf Facultad de Ciencias}\hspace{18mm}
\end{tabular}
\end{tabular}

\large
\vspace{10mm}
\hspace{25mm}
\begin{tabular}{l}
MÁSTER EN FÍSICA: RADIACIONES, NANOTECNOLOGÍA,\\
PARTÍCULAS Y ASTROFÍSICA

\end{tabular}


\vspace{5mm}
\hspace{25mm}
\begin{tabular}{l}
COMPLEMENTOS MATEMÁTICOS Y NUMÉRICOS
\\[1.5ex]
\LARGE\bf MONTE CARLO DE DIFUSIÓN Y \\
\LARGE\bf ÁTOMO ARMÓNICO
\end{tabular}
%

\hspace{2mm}
\begin{tabular}{l}
Por:
\\
{\bf A. S. Amari Rabah}\\


\end{tabular}
%




\begin{abstract}

Se estudia numéricamente la energía del estado fundamental en función del parámetro de interacción $\beta^2$ y el número de partículas $N$ de un sistema de partículas bosónicas bajo un potencial armónico mediante el método de Monte Carlo de difusión (DMC), considerando tanto la formulación básica sin \textit{importance sampling} como su versión con función guía usando la aproximación de Trotter.



\end{abstract}




% Indice %%%%%%%%%%%%%%%%%%%%%%%%%%%%%%%%%%%%%%%%%%%%%%%%%%%%%%%%%%%%%%%%%%%%%%%

\tableofcontents

% Texto %%%%%%%%%%%%%%%%%%%%%%%%%%%%%%%%%%%%%%%%%%%%%%%%%%%%%%%%%%%%%%%%%%%%%%%%

\pagestyle{fancy}
\fancyhead[RO,LE]{\leftmark}
\fancyhead[LO,RE]{\thepage}
\fancyfoot{}



\newpage


\section{Introducción}


Partiendo de la ecuación de Schrödinger dependiente del tiempo con una energía de referncia $E_T$

\begin{equation}
    i\hbar \frac{\partial \Psi (\mathbf{r}, t) }{\partial t} = (\hat{H} - E_T) \Psi (\mathbf{r}, t),
\end{equation}


y pasando a tiempo imaginario $\tau = i\hbar/t$ se obtiene 

\begin{equation}
    \frac{\partial \Psi (\mathbf{r}, \tau) }{\partial \tau} = -(\hat{H} - E_T) \Psi (\mathbf{r}, \tau)
\label{ESimag}
\end{equation}

donde $\mathbf{r}$ es el vector de posiciones de las particular en el espacio de configuraciones $\mathbf{r} \in \Omega, \mathbf{r} = (\mathbf{r}_1, \mathbf{r}_2, \cdots, \mathbf{r_N}).$ La ecuación \ref{ESimag} tiene solución \textit{simbólica}

$$
    \Psi (\mathbf{r}, \tau) = e^{-(\hat{H}-E_T)\tau} \Psi (\mathbf{r}, \tau=0).
$$


Suponiendo que la función de onda inicial se puede descomponer en autofunciones del hamiltoniano $\hat{H}\psi_n = E_n\psi_n$, con las energías $E_n$ crecientes ($E_1 < E_2 < \cdots < E_n$)

$$
    \Psi (\mathbf{r}, \tau) = \sum_{n} e^{-(E_n-E_T)\tau} \braket{\psi_n|\Psi_0} \psi_n.
$$

A tiempos $\tau$ grandes domina el estado fundamental, por lo que

$$
    \lim_{\tau \to \infty} \Psi(\mathbf{r}, \tau) \to e^{-(E_0-E_T)\tau} \braket{\psi_0|\Psi_0} \psi_0.
$$




El mayor inconveniente de aplicar el operador de evolución temporal para propagar el estado inicial en el tiempo es calcular la exponencial del hamiltoniano, ya que los operadores $\hat{T}$ y $\hat{V}$ no suelen conmutar

$$
    \ket{\Psi(\mathbf{r}, \tau)} = e^{-(\hat{H}-E_T)\tau} \ket{\Psi(\mathbf{r}, 0)}
$$

Por ello se aplica la \textit{aproximación de tiempo corto}, donde se calcula el estado a cada a paso de tiempo pequeño $\Delta \tau$

$$
    \ket{\Psi(\mathbf{r}, \tau + \Delta \tau)} = e^{-(\hat{H}-E_T)\Delta \tau} \ket{\Psi(\mathbf{r}, \tau)}.
$$

Aplicando la aproximación de Trotter al operador de evolución temporal \cite{kalos2008monte}

\begin{equation}
    e^{-(\hat{H}-E_T)\Delta \tau} = e^{-(\hat{V}-E_T)\Delta \tau/2}e^{-\hat{T}\Delta \tau}e^{-(\hat{V}-E_T)\Delta \tau /2} + \mathcal{O} (\Delta \tau^3).
\label{Trotter}
\end{equation}
    


Por lo que en representación de posiciones el estado a tiempo $\tau + \Delta \tau$ es

$$
    \Psi(\mathbf{r}, \tau + \Delta \tau) = \braket{\mathbf{r}|\Psi(\tau + \Delta \tau)} =  \braket{\mathbf{r}| e^{-(\hat{H}-E_T)\Delta \tau} | \Psi( \tau)}
$$

usando la relación de clausura en $\mathbf{r'}$

$$
    \Psi(\mathbf{r}, \tau + \Delta \tau) =\int_{\mathbf{r'}} \braket{\mathbf{r}| e^{-(\hat{H}-E_T)\Delta \tau} | \mathbf{r'}} \braket{{\mathbf{r'}} |\Psi( \tau )} d\mathbf{r'}
    = \int_{\mathbf{r'}} \braket{\mathbf{r}| e^{-(\hat{H}-E_T)\Delta \tau} | \mathbf{r'}} \Psi(\mathbf{r'}, \tau ) d\mathbf{r'}.
$$

Usando de nuevo la relación de clausura en $\mathbf{r''}$ y $\mathbf{r'''}$ con la aproximación \ref{Trotter}


$$
    \Psi(\mathbf{r}, \tau + \Delta \tau) =
    \int_{\mathbf{r'}} \int_{\mathbf{r''}} \int_{\mathbf{r'''}}
    \braket{\mathbf{r}| e^{-(\hat{V}-E_T)\Delta \tau/2} | \mathbf{r''}} 
$$
$$
    \cdot \braket{\mathbf{r''}| e^{-\hat{T}\Delta \tau} | \mathbf{r'''}}
    \braket{\mathbf{r'''}|e^{-(\hat{V}-E_T)\Delta \tau/2} | \mathbf{r'}}
    |\Psi(\mathbf{r'}, \tau ) 
    d\mathbf{r'} d\mathbf{r''} d\mathbf{r''}.
$$

Computando cada término individualmente

\begin{itemize}
    \item $\braket{\mathbf{r}| e^{-(\hat{V}-E_T)\Delta \tau/2} | \mathbf{r''}} =  e^{-(\hat{V}(\mathbf{r''})-E_T)} \braket{\mathbf{r|r''}} = e^{-(\hat{V}(\mathbf{r''})-E_T)} \delta(\mathbf{r-r''})$
    \item $ \braket{\mathbf{r'''}|e^{-(\hat{V}-E_T)\Delta \tau/2} | \mathbf{r'}} =  e^{-(\hat{V}(\mathbf{r'})-E_T)} \braket{\mathbf{r'''|r'}} = e^{-(\hat{V}(\mathbf{r'})-E_T)} \delta(\mathbf{r'''-r'})$
    \item El operador energía cinética es $-\Delta \tau \hat{T} = \frac{\Delta \tau \hbar^2}{2m} \nabla^2$, por lo que:
    
    $ \braket{\mathbf{r''}| e^{-\hat{T}\Delta \tau} | \mathbf{r'''}} = \frac{1}{(4\pi D \Delta \tau)^{N/2}} e^{-\frac{(\mathbf{r''} - \mathbf{r'''})^2}{4\pi D \Delta \tau}}$ 

    Donde $D=\hbar^2/2m$ y $\sigma^2 = 2D\Delta \tau$.
    
    
\end{itemize}

Por lo que la evolución total del estado es

\begin{equation}
    \Psi(\mathbf{r}, \tau + \Delta \tau) = \frac{1}{(2\pi \sigma^2)^{N/2}} \int_{\mathbf{r'}} e^{-\frac{(\mathbf{r} - \mathbf{r'}^2)}{2\sigma^2}} e^{-\Delta \tau /2 \left[ \hat{V}(\mathbf{r}) +  \hat{V}(\mathbf{r'}) - 2E_T) \right] } \Psi (\mathbf{r'}, \tau) d \mathbf{r'}.
\label{Psitotal}
\end{equation}




\subsection{Monte Carlo de difusión}
\label{DMC}


En la práctica, cuando el número de partículas $N$ es elevado la integral en \ref{Psitotal} es imposible de evaluar, por esta razón, se emplea una interpretación estocástica basada en métodos Monte Carlo, conocida como \textit{Diffusion Monte Carlo} (DMC) \cite{hammond1994monte}.



En el método DMC, la función de onda no se representa como una función continua, sino mediante un conjunto de $M$ partículas ficticias denominadas caminantes o \textit{walkers}. Cada walker corresponde a un punto en el espacio de configuraciones del sistema, es decir, tomamos $\Psi$ como

$$
    \Psi(\mathbf{r}, \tau) = \sum_{n=0}^M \delta(\mathbf{r} - \mathbf{r_i}(\tau))
$$

La densidad de walkers en el espacio es proporcional al valor de la función de onda. La ecuación de evolución tiene la forma general

\begin{center}
nuevo estado = propagador $\times$ estado anterior.
\end{center}

En el contexto Monte Carlo, el propagador se interpreta como una distribución de probabilidad que determina la transición de un walker desde una posición inicial hacia una nueva. Este propagador puede separarse en dos contribuciones principales: difusión y ramificación.\\


El término gaussiano del propagador describe un proceso de difusión. En cada paso temporal, cada walker se desplaza aleatoriamente siguiendo una distribución normal (la nueva posición se obtiene sumando un desplazamiento aleatorio gaussiano a la posición anterior). Este proceso representa el efecto de la energía cinética y puede interpretarse como una caminata aleatoria por el espacio de configuraciones.


El término dependiente del potencial no modifica directamente la posición del walker, sino su peso estadístico \footnote{En la simulación factor multiplicativo $m_i$ se toma como la parte entera del peso $w_i$ más un número aleatorio uniforme entre 0 y 1 para evitar un sesgo.}. Este factor determina la probabilidad de supervivencia o reproducción de cada walker. Cuando el potencial local es bajo, el walker tiene mayor probabilidad de reproducirse. En regiones de potencial elevado, los walkers tienden a desaparecer, es decir, este proceso actúa como un proceso de selección natural que favorece las regiones de baja energía\footnote{En de potenciales que tienden a $-\infty$ los caminantes se reproducen indefinidamente, por lo que estos casos requieren un tratamiento especial en DMC}.\\

De esta forma la ecuación de Langevin asociada a la Fokker--Planck \ref{ESimag} es \cite{van1992stochastic}

$$
    d \mathbf{r} (\tau) = \sqrt{2D} d \mathbf{W} (\tau),
$$

donde $d\mathbf{W}$ es un proceso de Weiner. En forma discreta el proceso es markoviano y la posición de los caminantes obedece la ecuación

\begin{equation}
    \mathbf{r_{n+1}} = \mathbf{r_n} + \sqrt{2D\Delta \tau}  \, \boldsymbol{\eta_n}
\label{langevindmcpuro}
\end{equation}
    
    

donde $\boldsymbol{\eta_n} \sim \mathcal{N}(0, 1)$ y se multiplican a cada paso de tiempo por el peso $w_n = \exp \left[ - \Delta \tau / 2 \cdot ( V(\mathbf{r_n}) + V(\mathbf{r_{n+1}})- 2E_T) \right]$.\\



\subsubsection{Estimación de la energía}



Para la estimación de la energía fundamental caben dos interpretaciones respecto a $E_T$. Si se integra la expresión \ref{Psitotal} sobre todas las posiciones $\mathbf{r}$ de los walkers se obtiene una estimación de la población $M(\tau)$

$$
    M(\tau + \Delta \tau) = \int_{\mathbf{r}} \Psi(\mathbf{r}, \tau + \Delta \tau) d\mathbf{r} = e^{-(E_0 - E_T)\Delta \tau} M(\tau)
$$


La primera opción es tomar una población de walkers que cambia con el tiempo (debido a los pesos) y $E_T$ constante, por lo que iterando en la relación anterior la energía del estado fundamental es

$$
    E_0 \approx E_T - \frac{1}{\tau} \ln \frac{M(\tau)}{M(0)}.
$$

O equivalentemente tomar el promedio de varias de medidas de

$$
    E_0 \approx E_T + \frac{1}{\Delta \tau} \ln \frac{M(\tau)}{M(\tau + \Delta \tau)}.
$$

\newpage
En este esquema la población al paso de tiempo siguiente se toma coma la población a tiempo anterior con cada caminante multiplicado por su peso estadístico, lo implica añadir algún método de control poblacional para evitar que el número de caminantes crezca demasiado. La otra opción es tomar una población de caminantes constante en el tiempo $M_0$ y $E_T (\tau)$ variable, es decir

$$
    \frac{d M(\tau)}{d\tau} = \int_{\mathbf{r}} \frac{\partial \Psi (\mathbf{r}, \tau)}{\partial \tau} d\mathbf{r} =  \int_{\mathbf{r}} (\hat{T} + \hat{V} - E_T (\tau)) \Psi (\mathbf{r}, \tau) d\mathbf{r} = \int_{\mathbf{r}} (E_0 - E_T(\tau)) \Psi (\mathbf{r}, \tau) d\mathbf{r} = 0\, .
$$


El término cinético $ \int_{\mathbf{r}} \nabla \cdot ( \nabla \Psi(\mathbf{r}, \tau) ) d\mathbf{r}$ se debe evaluar en los extremos del espacio de configuración, por lo que considerando que el gradiente de $\Psi(\mathbf{r}, \tau)$ decae lo suficientemente rápido esta contribución se puede despreciar. Particularmente, en este trabajo se considera un potencial de tipo armónico, por lo que esta suposición es válida.

La condición sobre $E_T$ es

$$
    E_T (\tau)= \frac{\int \hat{V} (\mathbf{r}) \Psi(\mathbf{r},  \tau) d\mathbf{r}}{\int  \Psi(\mathbf{r},  \tau) d\mathbf{r}} = \braket{\hat{V}}_{\Psi (\tau)} \, ,
$$

que se ajusta dinámicamente a cada paso de tiempo como

$$
    E_T (\tau + \Delta \tau)=  \braket{\hat{V}}_{\Psi (\tau)}.
$$

La energía del estado fundamental $E_0$ se estima como el promedio en el tiempo de $E_T(\tau)$. En este esquema la población de walkers a cada paso de tiempo se elige como una extracción uniforme de $M_0$ caminantes de la población a tiempo anterior, teniendo el cuenta el peso estadístico de cada uno debido al término de potencial. Esto podría introducir un sesgo si la población no es muy elevada, ya que los caminantes con pesos pequeños se ven apocados por aquellos con pesos mas elevados al restringir población a un numero fijo.\\

La primera opción presenta el inconveniente de introducir un método para controlar la población (ya que esta podría crecer demasiado y elevar el coste computacional). Mientras que el segundo restringe la población a un número fijo, lo que limita la representación de los caminantes en zonas elevadas del potencial.

Para evitar este problema se introduce un esquema mixto, es decir, se permite a la población $oscilar$ entre un valor objetivo $M_0$. De esta forma la energía de referencia a cada paso de tiempo se obtiene como

$$E_T(\tau + \Delta \tau) = \braket{\hat{V}}_{\Psi(\tau)} + \frac{\gamma}{\Delta \tau} \ln \left( \frac{M_0}{M(\tau)} \right),$$


donde $\gamma$ es un factor pequeño que evita que la oscilación sea muy elevada. La idea es que si la población $M(\tau)$ cae por debajo de $M_0$, el término se vuelve positivo, elevando $E_T$ y fomentando la creación de nuevos caminantes. Por el contrario, si la población excede el objetivo, $E_T$ disminuye para favorecer la aniquilación de caminantes.






\subsection{\textit{Importance Sampling}}
\label{IS}



Aunque el método de Monte Carlo de difusión arregla en gran medida la cuestión de las dimensiones, introduce otros problemas. Los walkers exploran el espacio de configuraciones mediante difusión libre visitando con alta probabilidad regiones donde la función de onda es prácticamente nula, por lo que gran parte del esfuerzo computacional se dedica a zonas irrelevantes para el estado fundamental. Lo que hace que la varianza estadística crezca rápidamente y las fluctuaciones llegan a ser muy elevadas.\\


Para mejorar la eficiencia, se introduce una función guía $\Psi_T(\mathbf r)$ que aproxima de alguna manera la forma del estado fundamental. En lugar de muestrear directamente $\Psi$, se define la función

\begin{equation}
f(\mathbf r,\tau)
=
\Psi_T(\mathbf r)\Psi(\mathbf r,\tau).
\end{equation}


Al sustituir esta definición en la ecuación de evolución \ref{ESimag}, se obtiene una ecuación de Fokker--Planck

\begin{equation}
\frac{\partial f}{\partial \tau}
=
D \nabla^2 f
-
\nabla \cdot \left( \mathbf F(\mathbf r) f \right)
-
\left[E_L(\mathbf r)-E_T\right] f,
\end{equation}

donde $\mathbf F(\mathbf r) = 2D \nabla \ln \Psi_T(\mathbf r)$ es la fuerza de deriva, y $E_L(\mathbf r) = \frac{\hat H \Psi_T(\mathbf r)}{\Psi_T(\mathbf r)}$ es la energía local. Ahora el término cinético se transforma en un proceso de difusión con deriva y el termino de ramificación del potencial depende solo de la energía local. La aproximación de Trotter para la evolución del estado queda como

\begin{equation}
\begin{split}
    f(\mathbf{r}, \tau + \Delta \tau) = \int_{\mathbf{r'}} \frac{1}{(2\pi \sigma^2)^{N/2}} \exp \left( -\frac{1}{2\sigma^2} \left[ \mathbf{r} - \mathbf{r'} -\mathbf{F}(\mathbf{r'}) \Delta \tau \right]^2 \right) \\
    \cdot \exp \left( \Delta \tau \left[ E_T - \frac{1}{2}( E_L(\mathbf{r}) + E_L(\mathbf{r'})) \right] \right) f(\mathbf{r'}, \tau) \, d\mathbf{r'}.
\end{split}
\end{equation}




Con esta transformación los caminantes ya no realizan un paseo aleatorio puro, sino que \textit{sienten} una deriva dirigida hacia las regiones donde $\Psi_T$ es grande. El muestreo se concentra en las zonas relevantes del espacio de configuraciones, la varianza estadística se reduce drásticamente y el proceso de ramificación depende de la energía local en lugar del potencial. La ecuación de Langevin asociada en versión discreta que dirige a los caminantes es

\begin{equation}
\mathbf r_{n+1}
=
\mathbf r_n
+
\mathbf F(\mathbf r_n)\Delta\tau
+
\sqrt{2D\Delta\tau}\,\boldsymbol\eta_n,
\notag
\end{equation}

que se multiplican por los pesos $w_n = \exp\left[-\Delta\tau / 2 \left(E_L(\mathbf r_{n}) + E_L(\mathbf r_{n+1}) - 2E_T\right)\right]$.


Finalmente, la energía del estado fundamental $E_0$ se estima como el valor medio de la energía local respecto a $f$, ya que usando la hermiticidad del operador hamiltoniano: 

$$
    \braket{E_L}_{f} = \frac{\int f E_L dr}{\int  f dr} = \frac{ \int \Psi_T \Psi_0 \frac{\hat{H} \Psi_T }{\Psi_T} dr}{ \int \Psi_T \Psi_0 } 
    = \frac{ \int \Psi_0 \hat{H} \Psi_T dr}{  \int \Psi_T \Psi_0 dr }
$$
$$
    = \frac{ \int \Psi_T \hat{H} \Psi_0 dr}{  \int \Psi_T \Psi_0 dr } = \frac{ \int \Psi_T E_0 \Psi_0 dr}{  \int \Psi_T \Psi_0 dr } = E_0
$$



El ajuste para mantener la población oscilando entorno a un objetivo $M_0$ en este caso es

$$E_T(\tau + \Delta \tau) = \braket{\hat{E_L}}_{f(\tau)} + \frac{\gamma}{\Delta \tau} \ln \left( \frac{M_0}{M(\tau)} \right).$$





\section{El átomo armónico}



Se consideran partículas que se mueven en un espacio tridimensional y que interactúan a través del hamiltoniano


\begin{equation}
    \hat{H} = -\frac{1}{2} \sum_{k=1}^N \nabla_{\mathbf{r}_k}^2 
    + \frac{1}{2} \sum_{k=1}^N \mathbf{r}_k^2  -\frac{\beta^2}{2N} \sum_{l>k=1}^N |\mathbf{r}_{k} - \mathbf{r}_l|^2 \,,
\label{hamiltonianoatomoarmonico}
\end{equation}


con $0 \leq \beta^2 <1$. La expresión analítica para la energía del estado fundamental es \cite{armstrong2012quantum}

\begin{equation}
    E_{exact} = \frac{3}{2} \left[ 1 + (N-1) \sqrt{(1 - \beta^2)}\right].
\label{Eexact}
\end{equation}
    





\subsection{Monte Carlo de difusión sin \textit{importance sampling}}




Puesto que el potencial es armónico los caminantes se inicializan mediante una distribución gaussiana centrada en el origen, $\mathbf r_i(0)\sim \mathcal N(0, D)$. Debido a que esta inicialización es arbitraria y no tiene por que representar el estado real del sistema los primeros pasos de tiempo se descartan, lo que se denomina termalización o \textit{burn-in} \cite{newman1999monte}. En este trabajo el tiempo de termalización se toma como \texttt{0.3 * Nsteps}, con $Nsteps$ el números de pasos de tiempo $\Delta \tau$.



En cada paso de tiempo $\Delta\tau$ los caminantes evolucionan en difusión libre a través de la ecuación de Langevin $\mathbf{r_{n+1}} = \mathbf{r_n} + \sqrt{2D\Delta \tau}  \, \boldsymbol{\eta_n}$ y se multiplican por el peso $w_n = \exp \left[ - \Delta \tau / 2 \cdot ( V(\mathbf{r_n}) + V(\mathbf{r_{n+1}})- 2E_T) \right]$.



El número total de caminantes se regula dinámicamente ajustando el parámetro $E_T$ a través de $E_T(\tau+\Delta\tau) = \braket{\hat{V}}_{\Psi(\tau)} + \frac{1}{\Delta\tau} \ln\frac{M_0}{M(\tau)}$, donde $M_0$ es la población objetivo. Una vez alcanzado el régimen estacionario el valor medio de $E_T$ proporciona una estimación de la energía del estado fundamental.




\subsubsection{Resultados}





En la figura \ref{E0DMCpuro} se muestran los resultados numéricos obtenidos mediante el método DMC sin \textit{importance sampling} descrito en la sección \ref{DMC} junto con la solución analítica dada por \ref{Eexact} para la energía del estado fundamental $E_0$ del hamiltoniano \ref{hamiltonianoatomoarmonico} para valores distintos del número de partículas.


En todos los casos la energía del estado fundamental decrece monótonamente al aumentar $\beta^2$, reflejando el debilitamiento del confinamiento debido al incremento de la interacción entre partículas. Las pequeñas desviaciones pueden atribuirse a las fluctuaciones estadísticas del muestreo y en menor medidad error asociado a la discretización temporal.



Para $N=2$, $N=3$ y $N=10$ se observa un acuerdo global entre los resultados numéricos y la solución analítica en todo el rango de $\beta^2$. Sin embargo, para $N=20$ se aprecia una diferencia más pronunciada, especialmente para valores pequeños e intermedios de $\beta^2$, donde el método DMC tiende a sobrestimar la energía del estado fundamental. El origen de esta desviación está en el aumemto del espacio de configuraciones con el número de partículas, lo que dificulta el muestreo eficiente mediante difusión libre. En este régimen, una fracción significativa de los caminantes explora regiones de baja probabilidad aumentando la varianza. Además, los efectos del control finito de la población y del paso temporal discreto se ven amplificados en sistemas de mayor tamaño.

Surge de esta manera la necesidad de incorporar el \textit{importance sampling} para obtener mejores resultados en sistemas con un número elevado de partículas o con interacciones más intensas.




\begin{figure}[H]
    \centering

    % Fila superior (0 grados)
    \begin{minipage}[t]{0.48\textwidth}
        \centering
        \includegraphics[width=\linewidth, height=5cm]{figures/E0DMCPuroN2.png} \\[2pt]
        \small (a) $N = 2$.
    \end{minipage}
    \hfill
    \begin{minipage}[t]{0.48\textwidth}
        \centering
        \includegraphics[width=\linewidth, height=5cm]{figures/E0DMCPuroN5.png} \\[2pt]
        \small (b) $N = 5$. 
    \end{minipage}

    \vspace{0.5cm}  % Espacio entre filas

    % Fila inferior (10 grados)
    \begin{minipage}[t]{0.48\textwidth}
        \centering
        \includegraphics[width=\linewidth, height=5cm]{figures/E0DMCPuroN10.png} \\[2pt]
        \small (c) $N = 10$.
    \end{minipage}
    \hfill
    \begin{minipage}[t]{0.48\textwidth}
        \centering
        \includegraphics[width=\linewidth, height=5cm]{figures/E0DMCPuroN20.png} \\[2pt]
        \small (d) $N = 20$.
    \end{minipage}

    \caption{Energía del estado fundamental en función de $\beta^2$ para distintos valores del número de partículas. En unidades atómicas, se toma una discretización $\Delta \tau = 0.001$, $10^4$ pasos de tiempo y una población objetivo de $2\cdot10^3$ caminantes. Tras inicializar el sistema la energía empieza a medirse después de $3\cdot10^3$ pasos de tiempo. Se han promediado los resultados sobre $10^3$ realizaciones.
    }
    \label{E0DMCpuro}
\end{figure}





\subsection{\textit{Importance sampling}}



Como función para el muestro de la importancia se toma

\begin{equation}
    \ket{\Psi_\alpha} = \left( \frac{\alpha}{\pi} \right)^{\frac{3}{4}N}
    \exp \left( - \frac{\alpha}{2} (\mathbf{r}^2_1 + \cdots +  \mathbf{r}^2_N  )\right),
\end{equation}


que coincide con el estado fundamental exacto de un oscilador armónico sin interacción con frecuencia $\omega=\alpha$. Además, lo caminantes se inicializan a través de esta función. La fuerza de deriva es

\begin{equation}
\mathbf F(\mathbf r)
=
2D \nabla \ln \Psi_\alpha(\mathbf r)
=
-2D\alpha
\left(
\mathbf r_1,\mathbf r_2,\ldots,\mathbf r_N
\right),
\end{equation}

lo que implica una fuerza restauradora dirigida hacia el origen del espacio de configuraciones. La ecuación de Langevin toma la forma

\begin{equation}
\mathbf r_{n+1}
=
\mathbf r_n
-
2D\alpha\,\mathbf r_n \Delta\tau
+
\sqrt{2D\Delta\tau}\,\boldsymbol\eta_n.
\end{equation}


Por otro lado, la energía local se obtiene como

\begin{equation}
\begin{split}
E_L(\mathbf r) = \frac{\hat H \Psi_T(\mathbf r)}{\Psi_T(\mathbf r)} = \frac{3}{2}N\alpha + \frac12(1-\alpha^2)\sum_{k=1}^N r_k^2 - \frac{\beta^2}{2N}\sum_{l>k} r_{kl}^2 \\
= \frac{3}{2}N\alpha + \frac12(1-\alpha^2 - \beta^2)\sum_{k=1}^N r_k^2 + \frac{\beta^2}{2N} \left( \sum_{k=1} r_{k} \right)^2.
\end{split}
\end{equation}






Con esta transformación el muestreo deja de ser uniforme en el espacio de configuraciones y pasa a estar concentrado en las regiones físicamente relevantes, ya que se espera que la función de onda del estado con interacción sea similar al estado con interacción.\\ 

So observa que la energía local depende del parámetro $\alpha$,  que controla la anchura de la gaussiana y determina el grado de localización de los caminantes en el espacio de configuraciones. En el caso ideal en que $\Psi_\alpha$ coincidiera con la función de onda exacta del estado fundamental, la energía local sería constante y la varianza se anularía. En el límite sin interacción ($\beta=0$), el valor exacto de $\alpha=1$ correspondiente al estado fundamental del oscilador armónico, para el cual se obtiene $E_L=\frac{3}{2}N$ constante. En presencia de interacción, el término proporcional a $\beta^2$ modifica la anchura efectiva de la función de onda, haciendo necesario ajustar el valor de $\alpha$. Una estimación de $\alpha$ puede obtenerse imponiendo que el coeficiente del término $\sum_k r_k^2$ en la energía local sea nulo, lo que lleva a la condición aproximada

\begin{equation}
\alpha_{\rm opt} \simeq \sqrt{1-\beta^2}.
\end{equation}


\newpage
\subsubsection{Resultados}


En la figura \ref{figures/E0DMCIS} se muestran los resultados numéricos obtenidos mediante el método DMC con \textit{importance sampling} descrito en la sección \ref{IS} junto con la solución analítica dada por \ref{Eexact} para la energía del estado fundamental $E_0$ del hamiltoniano \ref{hamiltonianoatomoarmonico} para valores distintos del número de partículas.\\


En este caso se observa que la coincidencia es total, incluso para valores elevados de $\beta^2$ y el número de partículas $N$. Se destaca especialmente el caso del número de partículas grande, ya que suponiendo una resolución equivalente con métodos deterministas muestreando $m$ puntos por dimensión el coste sería de $\mathcal{O}(m^{3N})$. Con un ejemplo modesto de 20 puntos por dimensión y 10 partículas se obtiene un coste de $20^{30} \sim 10^{39}$, aunque se realicen $10^{18}$ operaciones por segundo el tiempo restante sería mayor al del universo.





\begin{figure}[H]
    \centering


    \vspace{0.5cm}  % Espacio entre filas

    % Fila inferior (10 grados)
    \begin{minipage}[t]{0.48\textwidth}
        \centering
        \includegraphics[width=\linewidth, height=5cm]{figures/E0DMCISN10.png} \\[2pt]
        \small (c) $N = 10$.
    \end{minipage}
    \hfill
    \begin{minipage}[t]{0.48\textwidth}
        \centering
        \includegraphics[width=\linewidth, height=5cm]{figures/E0DMCISN20.png} \\[2pt]
        \small (d) $N = 20$.
    \end{minipage}

    % Fila inferior (10 grados)
    \begin{minipage}[t]{0.48\textwidth}
        \centering
        \includegraphics[width=\linewidth, height=5cm]{figures/E0DMCISN50.png} \\[2pt]
        \small (c) $N = 50$.
    \end{minipage}
    \hfill
    \begin{minipage}[t]{0.48\textwidth}
        \centering
        \includegraphics[width=\linewidth, height=5cm]{figures/E0DMCISN100.png} \\[2pt]
        \small (d) $N = 100$.
    \end{minipage}

    \caption{Energía del estado fundamental en función de $\beta^2$ para distintos valores del número de partículas. En unidades atómicas, se toma una discretización $\Delta \tau = 0.001$, $10^4$ pasos de tiempo y una población objetivo de $2\cdot10^3$ caminantes. Tras inicializar el sistema la energía empieza a medirse después de $3\cdot10^3$ pasos de tiempo. Se han promediado los resultados sobre $10^3$ realizaciones.
    }
    \label{figures/E0DMCIS}
\end{figure}

















\newpage
\section{Conclusión}



Se ha estudiado numéricamente la energía del estado fundamental en función del parámetro de interacción $\beta^2$ y el número de partículas $N$ de un sistema de partículas bosónicas bajo un potencial armónico mediante el método de Monte Carlo de difusión (DMC), considerando tanto la formulación básica sin \textit{importance sampling} como su versión con función guía usando la aproximación de Trotter.

El método DMC puro muestra un buen acuerdo con la solución exacta para sistemas con un número reducido de partículas. Sin embargo, al aumentar el número de partículas, el DMC sin \textit{importance sampling} presenta un incremento notable en el error, debid0 a la exploración ineficiente del espacio de configuraciones.

Usando \textit{importance sampling} mediante una función guía gaussiana se añade un término de deriva que dirige a los caminantes hacia las regiones más relevantes del espacio de configuraciones y reemplaza el potencial por la energía local en el proceso de ramificación. Los resultados obtenidos con IS-DMC muestran un excelente acuerdo con la solución analítica en todo el rango de valores de $\beta^2$ considerado, incluso para sistemas con un número elevado de partículas.









\newpage


\bibliographystyle{plain}
\bibliography{bibliografiaDMC}





\end{document}